% GENERATED FILE. Edit canonical lessons, then run npm run generate:books.
% canonical-source-hash: fnv1a64:8cc6adcb90ffadba
% canonical-lessons: AR-C406-tatallaba, AR-C406-tawaffara, AR-C406-kafa, AR-C406-iftaqara, AR-C406-ijra

\chapter{What it takes}
\label{ch:ar-what-it-takes}
\hlchaptermodality{406}

\begin{chapteropening}
\textbf{By the end of this chapter:} \emph{I can say that something requires, is available, suffices or is lacking, and talk about a procedure.}

Complete the last lesson of chapter 406: I can say that something requires, is available, suffices or is lacking, and talk about a procedure.
\end{chapteropening}

\section[\texorpdfstring{\ar{تطلب} (taṭallaba)}{تطلب (taṭallaba)}]{\ar{تطلب} (taṭallaba) — to require}

\label{lesson:AR-C406-tatallaba}



\begin{quote}
Before the new one: say the Arabic for to contain, then the Arabic for to concern.
\end{quote}

\subsection*{You'll want to know: \ar{تطلب}}
\textbf{\ar{تطلب}} — \emph{taṭallaba} — \textquotedblleft{}to require\textquotedblright{}.

\begin{quote}
\ar{يتطلب} \ar{هذا} \ar{العمل} \ar{صبرا} \ar{كبيرا}.
\end{quote}

\begin{quote}
\emph{yataṭallabu hādhā al-ʿamal ṣabran kabīran.}
\end{quote}

\begin{quote}
This work requires great patience.
\end{quote}

Present: \ar{يتطلب}, \emph{yataṭallab}.

\begin{cousinweb}
\textbf{\ar{تطلب}} (\emph{taṭallaba}) is built on the native Arabic root \emph{ṭ-l-b}, \textquotedblleft{}to seek, to ask for\textquotedblright{} (Lane). It is Form V, to seek a thing again and again: to call for, to require.
\end{cousinweb}

\begin{grammarlens}[title={What you've built}]
One more word for saying what happens if.
\end{grammarlens}

\subsection*{Your turn}
\begin{itemize}
\raggedright
  \item \emph{Say it:} \emph{taṭallaba}
  \item \emph{Say it:} \emph{taṭallaba}, once more
  \item \emph{Say it:} \emph{yataṭallabu hādhā al-ʿamal ṣabran kabīran.} — This work requires great patience.
  \item \emph{Recall:} say the Arabic for to contain, then the Arabic for to concern, then say \emph{taṭallaba} again
\end{itemize}

\begin{tcolorbox}[breakable,colback=teal!4,colframe=teal!35!black,title={Before you move on}]
What does \ar{تطلب} mean? (\textquotedblleft{}To require\textquotedblright{}.) Say it once more.
\end{tcolorbox}
\section[\texorpdfstring{\ar{توفر} (tawaffara)}{توفر (tawaffara)}]{\ar{توفر} (tawaffara) — to be available}

\label{lesson:AR-C406-tawaffara}



\begin{quote}
Before the new one: say the Arabic for to concern, then the Arabic for to require.
\end{quote}

\subsection*{You'll want to know: \ar{توفر}}
\textbf{\ar{توفر}} — \emph{tawaffara} — \textquotedblleft{}to be available\textquotedblright{}.

\begin{quote}
\ar{تتوفر} \ar{الغرف} \ar{في} \ar{الفندق} \ar{هذا} \ar{الأسبوع}.
\end{quote}

\begin{quote}
\emph{tatawaffaru al-ghuraf fī al-funduq hādhā al-usbūʿ.}
\end{quote}

\begin{quote}
Rooms are available at the hotel this week.
\end{quote}

Present: \ar{يتوفر}, \emph{yatawaffar}.

\begin{cousinweb}
\textbf{\ar{توفر}} (\emph{tawaffara}) is built on the native Arabic root \emph{w-f-r}, \textquotedblleft{}to be abundant, ample\textquotedblright{} (Lane). It is Form V, to become ample: to be available. \ar{وفر}, to save, is from the same root.
\end{cousinweb}

\begin{grammarlens}[title={What you've built}]
One more word for saying what happens if.
\end{grammarlens}

\subsection*{Your turn}
\begin{itemize}
\raggedright
  \item \emph{Say it:} \emph{tawaffara}
  \item \emph{Say it:} \emph{tawaffara}, once more
  \item \emph{Say it:} \emph{tatawaffaru al-ghuraf fī al-funduq hādhā al-usbūʿ.} — Rooms are available at the hotel this week.
  \item \emph{Recall:} say the Arabic for to concern, then the Arabic for to require, then say \emph{tawaffara} again
\end{itemize}

\begin{tcolorbox}[breakable,colback=teal!4,colframe=teal!35!black,title={Before you move on}]
What does \ar{توفر} mean? (\textquotedblleft{}To be available\textquotedblright{}.) Say it once more.
\end{tcolorbox}
\section[\texorpdfstring{\ar{كفى} (kafā)}{كفى (kafā)}]{\ar{كفى} (kafā) — to suffice}

\label{lesson:AR-C406-kafa}



\begin{quote}
Before the new one: say the Arabic for to require, then the Arabic for to be available.
\end{quote}

\subsection*{You'll want to know: \ar{كفى}}
\textbf{\ar{كفى}} — \emph{kafā} — \textquotedblleft{}to suffice\textquotedblright{}.

\begin{quote}
\ar{يكفي} \ar{هذا} \ar{الطعام} \ar{لثلاثة} \ar{أشخاص}.
\end{quote}

\begin{quote}
\emph{yakfī hādhā aṭ-ṭaʿām li-thalāthat ashkhāṣ.}
\end{quote}

\begin{quote}
This food is enough for three people.
\end{quote}

Present: \ar{يكفي}, \emph{yakfī}.

\begin{cousinweb}
\textbf{\ar{كفى}} (\emph{kafā}) is built on the native Arabic root \emph{k-f-y}, usually traced to the sense \textquotedblleft{}to be enough\textquotedblright{}. \emph{Kafā} is a Form I verb with a weak last letter: to be enough. Its present is \emph{yakfī}.
\end{cousinweb}

\begin{grammarlens}[title={What you've built}]
One more word for saying what happens if.
\end{grammarlens}

\subsection*{Your turn}
\begin{itemize}
\raggedright
  \item \emph{Say it:} \emph{kafā}
  \item \emph{Say it:} \emph{kafā}, once more
  \item \emph{Say it:} \emph{yakfī hādhā aṭ-ṭaʿām li-thalāthat ashkhāṣ.} — This food is enough for three people.
  \item \emph{Recall:} say the Arabic for to require, then the Arabic for to be available, then say \emph{kafā} again
\end{itemize}

\begin{tcolorbox}[breakable,colback=teal!4,colframe=teal!35!black,title={Before you move on}]
What does \ar{كفى} mean? (\textquotedblleft{}To suffice\textquotedblright{}.) Say it once more.
\end{tcolorbox}
\section[\texorpdfstring{\ar{افتقر} (iftaqara)}{افتقر (iftaqara)}]{\ar{افتقر} (iftaqara) — to lack}

\label{lesson:AR-C406-iftaqara}



\begin{quote}
Before the new one: say the Arabic for to be available, then the Arabic for to suffice.
\end{quote}

\subsection*{You'll want to know: \ar{افتقر}}
\textbf{\ar{افتقر}} — \emph{iftaqara} — \textquotedblleft{}to lack\textquotedblright{}.

\begin{quote}
\ar{تفتقر} \ar{القرية} \ar{إلى} \ar{طبيب}.
\end{quote}

\begin{quote}
\emph{taftaqiru al-qarya ilā ṭabīb.}
\end{quote}

\begin{quote}
The village lacks a doctor.
\end{quote}

Present: \ar{يفتقر}, \emph{yaftaqir}.

It takes \ar{إلى}: \ar{افتقر} \ar{إلى} ….

\begin{cousinweb}
\textbf{\ar{افتقر}} (\emph{iftaqara}) is built on the native Arabic root \emph{f-q-r}, \textquotedblleft{}to be needy, poor\textquotedblright{} (Lane). It is Form VIII, to become needy of a thing: to lack it.
\end{cousinweb}

\begin{grammarlens}[title={What you've built}]
One more word for saying what happens if.
\end{grammarlens}

\subsection*{Your turn}
\begin{itemize}
\raggedright
  \item \emph{Say it:} \emph{iftaqara}
  \item \emph{Say it:} \emph{iftaqara}, once more
  \item \emph{Say it:} \emph{taftaqiru al-qarya ilā ṭabīb.} — The village lacks a doctor.
  \item \emph{Recall:} say the Arabic for to be available, then the Arabic for to suffice, then say \emph{iftaqara} again
\end{itemize}

\begin{tcolorbox}[breakable,colback=teal!4,colframe=teal!35!black,title={Before you move on}]
What does \ar{افتقر} mean? (\textquotedblleft{}To lack\textquotedblright{}.) Say it once more.
\end{tcolorbox}
\section[\texorpdfstring{\ar{إجراء} (ijrāʾ)}{إجراء (ijrāʾ)}]{\ar{إجراء} (ijrāʾ) — a procedure}

\label{lesson:AR-C406-ijra}



\begin{quote}
Before the new one: say the Arabic for to suffice, then the Arabic for to lack.
\end{quote}

\subsection*{You'll want to know: \ar{إجراء}}
\textbf{\ar{إجراء}} — \emph{ijrāʾ} — \textquotedblleft{}a procedure\textquotedblright{}.

\begin{quote}
\ar{الإجراء} \ar{بسيط} \ar{ويأخذ} \ar{دقيقة}.
\end{quote}

\begin{quote}
\emph{al-ijrāʾ basīṭ wa-yaʾkhudhu daqīqa.}
\end{quote}

\begin{quote}
The procedure is simple and takes a minute.
\end{quote}

\begin{cousinweb}
\textbf{\ar{إجراء}} (\emph{ijrāʾ}) is built on the native Arabic root \emph{j-r-y}, \textquotedblleft{}to run, to flow\textquotedblright{} (Lane). \emph{Ijrāʾ} is the Form IV verbal noun, making a thing run: carrying out, a procedure.
\end{cousinweb}

\begin{grammarlens}[title={What you've built}]
One more word for saying what happens if.
\end{grammarlens}

\subsection*{Your turn}
\begin{itemize}
\raggedright
  \item \emph{Say it:} \emph{ijrāʾ}
  \item \emph{Say it:} \emph{ijrāʾ}, once more
  \item \emph{Say it:} \emph{al-ijrāʾ basīṭ wa-yaʾkhudhu daqīqa.} — The procedure is simple and takes a minute.
  \item \emph{Recall:} say the Arabic for to suffice, then the Arabic for to lack, then say \emph{ijrāʾ} again
\end{itemize}

\begin{tcolorbox}[breakable,colback=teal!4,colframe=teal!35!black,title={Before you move on}]
What does \ar{إجراء} mean? (\textquotedblleft{}A procedure\textquotedblright{}.) Say it once more.
\end{tcolorbox}
